\section{Evolución semi-implícita estabilizada de la frontera libre}
\label{sec:semiimplicito}

\subsection{Diagnóstico de la restricción del paso de tiempo}

El método de Características de Galerkin (\ref{eqn:nuevophih}) es incondicionalmente estable para la advección lineal $\phi_t+S\cdot\nabla\phi=0$ con $S$ dado, de modo que la restricción sobre $\Delta t$ no proviene del transporte sino del \emph{acoplamiento explícito} entre la forma de la frontera y su velocidad. En el modelo adimensional, la velocidad normal es
\[
V_n=-\partial_n P+G\,\partial_n U-AG\,\frac{n\cdot x}{2},\qquad \Delta P=0\ \text{en }\Omega,\qquad P|_\Gamma=\kappa-AG\,\frac{|x|^2}{4}.
\]
Si $\Gamma$ es una perturbación de una circunferencia de radio $R$ con modo de Fourier $k$ y amplitud $\eta$, entonces $\delta\kappa\sim k^2\eta/R^2$ y, al resolver el problema armónico para $P$, $\partial_n\delta P\sim (k/R)\,\delta\kappa$. La dinámica linealizada es
\[
\eta_t=-\lambda_k\,\eta,\qquad \lambda_k\sim \frac{k^3}{R^3}\quad (k\gg 1),
\]
es decir, el problema de frontera libre es de \emph{tercer orden} en el número de onda (tipo Hele--Shaw con tensión superficial). Con una malla de paso $h$ el modo más alto representable es $k_{\max}\sim\pi/h$, y cualquier esquema explícito en la curvatura exige
\[
\Delta t\lesssim \frac{2}{\lambda_{k_{\max}}}\sim C\,h^{3}.
\]
Esto explica la inestabilidad observada y el $\Delta t$ muy pequeño. Un segundo efecto adverso del método de características con interpolación $\mathbb{P}_1$ es la difusión numérica, de orden $h^2/\Delta t$ por unidad de tiempo, que \emph{crece} al disminuir $\Delta t$ y suaviza artificialmente la frontera.

\subsection{Esquema propuesto}

Se propone sustituir el paso (\ref{eqn:nuevophih}) por un esquema de elementos finitos continuos en $\Phi_h$ con tres ingredientes.

\paragraph{(i) Estabilizador semi-implícito (tipo Smereka / Hou--Lowengrub--Shelley).}
Se suma y se resta un operador implícito que reproduce el símbolo dominante $\lambda_k$ con un operador local de segundo orden con coeficiente proporcional a $1/h$. Sea $\delta\phi=\phi^{n+1}-\phi^n$. Como $|\nabla\phi^n|\approx 1$ (función distancia), el término estabilizador es $\theta_h\Delta\,\delta\phi$ con
\begin{equation}
	\theta_h = c_\theta\,\frac{\gamma_{ref}}{h},\qquad c_\theta\ge \frac{\pi}{2},
	\label{eqn:thetah}
\end{equation}
donde $\gamma_{ref}$ es la constante que multiplica a $\kappa$ en el modelo (en las variables adimensionales $\gamma_{ref}=1$). El término se anula en el límite $\Delta t\to 0$, de modo que el esquema es consistente de orden uno en el tiempo.

\paragraph{(ii) Transporte con elementos finitos continuos estabilizados.}
En lugar de interpolar $\phi^n$ en el pie de la característica, se resuelve la forma débil con estabilización SUPG o de penalización de saltos del gradiente (CIP), sin difusión de interpolación. Como el término estabilizador $\theta_h$ ya cubre la rigidez, el transporte puede tratarse con $\theta$-esquema (Crank--Nicolson) o SSP-RK2.

\paragraph{(iii) Una sola extensión escalar.}
Como $\partial_t\phi+V_n|\nabla\phi|=0$, basta extender la velocidad normal escalar $F$ (un solo problema elíptico (\ref{eqn:problemaSh}) en lugar de dos para $S_1,S_2$) y definir $S^n=F^n\,\nabla\phi^n/|\nabla\phi^n|$.

\paragraph{Problema discreto.} Dados $\phi^n$ y $S^n$ (calculados con la curvatura, $c_h$, $p_h$ y la extensión en $t_n$), encontrar $\phi^{n+1}\in\Phi_h$ tal que, para todo $v_h\in\Phi_h$,
\begin{equation}
	\begin{aligned}
		\Big(\frac{\phi^{n+1}-\phi^n}{\Delta t},\,v_h+\tau_K S^n\!\cdot\!\nabla v_h\Big)_{\mathcal{O}}
		&+\theta_h\big(\nabla(\phi^{n+1}-\phi^n),\nabla v_h\big)_{\mathcal{O}}\\
		&+\Big(S^n\!\cdot\!\nabla\phi^{n+1/2},\,v_h+\tau_K S^n\!\cdot\!\nabla v_h\Big)_{\mathcal{O}}=0,
	\end{aligned}
	\label{eqn:phisemiimplicito}
\end{equation}
con $\phi^{n+1/2}=\tfrac12(\phi^{n}+\phi^{n+1})$ y parámetro SUPG $\tau_K=\dfrac{h_K}{2\,\|S^n\|_{\infty,K}}$ (valor estándar de SUPG; se puede ajustar experimentalmente). La matriz $\tfrac1{\Delta t}M+\tfrac12 C(S^n)+\theta_h K$ es dispersa y se ensambla una vez por paso; el costo es una resolución adicional comparable a las ya presentes (curvatura, $c_h$, $p_h$, extensión).

\subsection{Estabilidad del estabilizador}

Para el modo $k$ del problema linealizado, la ecuación $\eta_t=-\lambda_k\eta$ con $\lambda_k\le\gamma_{ref}\,k^3$ discretizada con el estabilizador se reduce a
\[
(1+\Delta t\,\theta_h k^2)(\eta^{n+1}-\eta^n)=-\Delta t\,\lambda_k\,\eta^n
\ \Longrightarrow\
g_k=1-\frac{\Delta t\,\lambda_k}{1+\Delta t\,\theta_h k^2}.
\]
Entonces $|g_k|\le 1$ equivale a $\dfrac{\Delta t\,\lambda_k}{1+\Delta t\,\theta_h k^2}\le 2$, que se cumple para todo $\Delta t>0$ siempre que $\theta_h k^2\ge \lambda_k/2$, es decir $\theta_h\ge \gamma_{ref}\,k_{\max}/2=\gamma_{ref}\pi/(2h)$, que es (\ref{eqn:thetah}). Así, la restricción $\Delta t\lesssim h^3$ desaparece y $\Delta t$ queda limitado únicamente por la precisión y por el desplazamiento de la frontera por paso ($\max|V_n|\Delta t\lesssim h$).

Verificación numérica del modelo lineal (con $\gamma_{ref}=1$, $\theta_h=2/h$, $\Delta t=h$, modos $1\le k\le\pi/h$):
\begin{center}
\begin{tabular}{ccccc}
\toprule
$h$ & $\Delta t_{\max}$ explícito $\approx 2/k_{\max}^3$ & $\Delta t$ semi-implícito & $\max_k|g_k|$ & razón\\
\midrule
$0.100$ & $6.5\times10^{-5}$ & $0.100$ & $0.967$ & $\approx 1.6\times10^{3}$\\
$0.050$ & $8.1\times10^{-6}$ & $0.050$ & $0.983$ & $\approx 6.2\times10^{3}$\\
$0.025$ & $1.0\times10^{-6}$ & $0.025$ & $0.992$ & $\approx 2.5\times10^{4}$\\
\bottomrule
\end{tabular}
\end{center}
Estos números corresponden al modelo lineal de Fourier, no al código completo; la ganancia real debe confirmarse con las simulaciones del capítulo de resultados numéricos.

\subsection{Algoritmo modificado}

Se reemplaza el último paso del Algoritmo~\ref{eqn:algoritmo}:
\begin{enumerate}
	\item Calcular $\kappa_h$, $c_h$, $p_h$ y $V_n$ en $t_n$ (como en el Algoritmo~\ref{eqn:algoritmo}).
	\item Extender $V_n$ a $F_h\in M_h$ resolviendo (\ref{eqn:problemaSh}) una sola vez y formar $S^n=F_h\,\nabla\phi_h^n/|\nabla\phi_h^n|$.
	\item Resolver (\ref{eqn:phisemiimplicito}) para $\phi_h^{n+1}$.
	\item Reinicializar $\phi_h^{n+1}$ \emph{solo si} $\max_{\Omega_h^\Gamma}\big||\nabla\phi_h^{n+1}|-1\big|>\mathrm{tol}$ (la extensión por $\nabla\phi\cdot\nabla u=0$ conserva la propiedad de distancia, por lo que no hace falta reinicializar en cada paso).
	\item Elegir $\Delta t_{n+1}=\mathrm{CFL}\;h/\max|F_h|$ con $\mathrm{CFL}\in[0.5,1]$.
\end{enumerate}

\subsection{Variante de mayor robustez (opcional)}

Si $\Delta t\sim h$ resultara aún insuficiente por la precisión de primer orden del estabilizador, se puede tratar la curvatura de forma totalmente implícita: linealizar $\kappa(\phi^{n+1})\approx\kappa^n-\Delta_\Gamma\eta-|II|^2\eta$ con $\eta=-\Delta t\,V_n^{n+1}$ y resolver el sistema acoplado $(P,\eta)$ por iteración de punto fijo o Newton--Krylov, usando el esquema anterior como precondicionador. Es más costoso por paso, pero admite $\Delta t$ aún mayores.

\subsection{Plan de validación}
\begin{enumerate}
	\item Solución radial exacta (Sección ``Solución exacta del crecimiento tumoral con simetría radial''): orden de convergencia en $\Delta t$ y $h$.
	\item Circunferencia perturbada con un modo $k$: comparar la tasa de decaimiento/crecimiento con la relación de dispersión lineal y medir el $\Delta t_{\max}$ estable frente a $h$ (esperado $\sim h^3$ con características, $\sim h$ con el nuevo esquema).
	\item Comparar tiempo total de CPU y conservación del área/volumen frente al método de características para el mismo tiempo final.
\end{enumerate}

\subsection{Verificación en el código FreeFem++}

El esquema se implementó en \texttt{edp/crecimientotumoranocircularSI.edp}, que reemplaza únicamente la llamada a \texttt{convect} del programa \texttt{crecimientotumoranocircularH.edp}. Se comparó con el método de características sobre el experimento A ($r_0(\theta)=2.1+0.5\cos 2\theta$, $A=0.5$, $G=20$), malla $60\times60$ ($h\approx0.19$), tiempo final $t=0.12$. La referencia es el método de características con $\Delta t=0.001$ y el error es la máxima diferencia de la posición radial de la frontera en 12 ángulos.

\begin{center}
\begin{tabular}{ccc}
\toprule
$\Delta t$ & Características & Semi-implícito ($c_\theta=2$)\\
\midrule
$0.010$ & $0.008$ & $0.027$\\
$0.020$ & $0.067$ & $0.025$\\
$0.030$ & $0.345$ & $0.026$\\
$0.040$ & $0.158$ & $0.021$\\
\bottomrule
\end{tabular}
\end{center}

Resultados y limitaciones:
\begin{itemize}
	\item El método de características pierde precisión a partir de $\Delta t\approx0.02$ y es inutilizable en $0.03$; el esquema propuesto mantiene un error casi constante hasta $\Delta t=0.04$. En esta malla y este experimento la ganancia observada es de un factor 2--4, \emph{no} los órdenes de magnitud del modelo lineal; la ganancia debe crecer al refinar la malla, pero eso no se ha medido todavía.
	\item El esquema propuesto introduce un sesgo (error de $0.01$--$0.03$) que no disminuye con $\Delta t$ y que depende de $c_\theta$: a $t=0.05$ y $\Delta t=0.005$ el error fue $0.013, 0.010, 0.007, 0.006$ para $c_\theta=2, 1, 0.5, 0.25$. Por ello el programa usa $c_\theta=0.5$ por defecto. La cota teórica $c_\theta\ge\pi/2$ es conservadora; falta determinar experimentalmente cuán pequeño puede ser $c_\theta$ antes de perder estabilidad con mallas más finas.
	\item Para $\Delta t$ pequeño ($\le 0.005$) el método de características es más preciso; la ventaja del nuevo esquema está en los pasos grandes.
	\item Solo se probó una geometría, una malla y tiempos cortos; no se observó explosión (\emph{blow-up}) en ninguno de los métodos, solo pérdida de precisión.
\end{itemize}

\paragraph{Mallas más finas.} Se repitió la comparación con mallas $100\times100$ ($h\approx0.113$) y $150\times150$ ($h\approx0.075$), $t=0.04$, referencia: características con $\Delta t=0.002$. Error máximo en la posición radial de la frontera:

\begin{center}
\begin{tabular}{c|ccc|ccc}
\toprule
 & \multicolumn{3}{c|}{$n=100$} & \multicolumn{3}{c}{$n=150$}\\
$\Delta t$ & Caract. & SI $c_\theta=0.5$ & SI $c_\theta=1$ & Caract. & SI $c_\theta=0.5$ & SI $c_\theta=1$\\
\midrule
$0.004$ & $0.0009$ & $0.0084$ & $0.0105$ & $0.0007$ & $0.0087$ & $0.0114$\\
$0.008$ & $0.0015$ & $0.0092$ & $0.0123$ & $0.0292$ & $0.0099$ & $0.0134$\\
$0.020$ & $0.0232$ & $0.0112$ & $0.0162$ & $0.0269$ & $0.0106$ & $0.0156$\\
$0.040$ & $0.0131$ & $0.0103$ & $0.0155$ & $0.0194$ & $0.0098$ & $0.0150$\\
\bottomrule
\end{tabular}
\end{center}

El método de características pierde precisión a partir de $\Delta t\approx0.02$ para $n=100$ y ya en $\Delta t=0.008$ para $n=150$ (el $\Delta t$ útil disminuye al refinar, como predice el análisis), mientras que el esquema semi-implícito mantiene un error casi constante ($\approx0.01$) hasta $\Delta t=0.04$ en todas las mallas. En $n=150$ la ventaja medida es de un factor $\ge5$ en el paso de tiempo útil (de $0.004$ a $\ge0.04$, que fue el mayor probado, de modo que es una cota inferior). No se observó \emph{blow-up} en ninguna corrida. El sesgo del esquema semi-implícito ($\approx0.01$) no disminuye al refinar la malla ni al reducir $\Delta t$; su origen no está aislado todavía, y por ahora limita la precisión alcanzable con este esquema frente al método de características con $\Delta t$ pequeño.

\paragraph{Origen del sesgo (prueba con $c_\theta=0$).} Con $n=100$, $t=0.04$ y la misma referencia, se repitió el esquema sin estabilizador ($c_\theta=0$), con y sin SUPG:

\begin{center}
\begin{tabular}{cccc}
\toprule
$\Delta t$ & $c_\theta=0$, SUPG & $c_\theta=0$, sin SUPG & $c_\theta=0.5$, SUPG\\
\midrule
$0.002$ & $0.0019$ & $0.0019$ & ---\\
$0.004$ & $0.0009$ & $0.0021$ & $0.0084$\\
$0.008$ & $0.0018$ & $0.0028$ & $0.0092$\\
$0.020$ & $0.0049$ & $0.0063$ & $0.0112$\\
\bottomrule
\end{tabular}
\end{center}

Conclusiones: (i) el sesgo de $\approx0.01$ proviene del estabilizador $\theta_h$ y no del transporte (Crank--Nicolson, SUPG, $\mathbb{P}_2$), que da errores de $0.001$--$0.005$; (ii) SUPG aporta una mejora modesta; (iii) con $\theta_h=0$ el esquema es estable incluso con $\Delta t=0.02$ en este experimento, es decir, en el experimento A (frontera suave, solo el modo $k=2$) el estabilizador \emph{no es necesario} y estos ensayos no demuestran todavía su efecto estabilizante. El error se concentra en los puntos de menor radio ($\theta=\pi/2,3\pi/2$), donde la frontera queda retrasada hasta $0.008$ respecto de la referencia. El sesgo tampoco crece linealmente con $c_\theta$ ($0.006$ con $0.25$ y $0.013$ con $2$ a $n=60$), por lo que su mecanismo no está identificado. Además, la afirmación anterior de que el término estabilizador se anula con $\Delta t\to0$ en el sentido de dar un error $O(\Delta t)$ no se observó en los datos (el error casi no cambia entre $\Delta t=0.004$ y $0.04$). Quedan por probar: (a) un caso con frontera irregular y modos altos (experimentos B1 o C), que es donde el método de características exige $\Delta t$ muy pequeño y donde el estabilizador debería ser necesario; (b) limitar el estabilizador a una banda alrededor de $\Gamma$.

\paragraph{Experimento B1 (frontera irregular, modos hasta $k=6$).} Malla $60\times60$ sobre $[-6,6]^2$, $t=0.04$, referencia: características con $\Delta t=0.00025$; error máximo en la posición radial en 48 ángulos.

\begin{center}
\begin{tabular}{c|c|cccc}
\toprule
$\Delta t$ & Caract. & SI $c_\theta=0$ & SI $c_\theta=0.25$ & SI $c_\theta=0.5$ & SI $c_\theta=1$\\
\midrule
$0.002$ & $0.0018$ & $0.0113$ & $0.0134$ & $0.0169$ & $0.0201$\\
$0.004$ & $0.0040$ & $0.0091$ & $0.0154$ & $0.0185$ & $0.0222$\\
$0.008$ & $0.0079$ & --- & --- & --- & ---\\
$0.010$ & --- & $0.0093$ & $0.0178$ & $0.0224$ & $0.0279$\\
$0.020$ & --- & $0.0118$ & $0.0147$ & $0.0235$ & $0.0377$\\
\bottomrule
\end{tabular}
\end{center}

En este horizonte temporal corto \emph{no} se reproduce la inestabilidad: el método de características es estable y su error crece linealmente con $\Delta t$ (para $\Delta t=0.008$ es $0.008$), sin ninguna corrida con \emph{blow-up}, y es más preciso que el esquema semi-implícito para todo $\Delta t$ probado. El esquema propuesto no aporta ventaja en este ensayo. La inestabilidad que obliga a usar $\Delta t=0.00025$ en B1 (tiempo final $t=1$) debe aparecer a tiempos mayores, cuando la frontera desarrolla curvatura alta; no está probado todavía que el esquema propuesto la evite.

\subsection{Comparación con características de Galerkin (proyección $L^2$)}

El programa \texttt{edp/crecimientotumoranocircularGalerkin.edp} sustituye la interpolación nodal \texttt{phih = convect(\dots)} por la proyección $L^2$ $(\phi^{n+1},v)=(\phi^n\circ X,v)$, con $\phi_h\in\mathbb{P}_2$ y $S\in\mathbb{P}_1$.

\paragraph{Experimento A ($n=100$, $t=0.04$).} Error máximo radial frente a la referencia nodal con $\Delta t=0.002$:
\begin{center}
\begin{tabular}{ccc}
\toprule
$\Delta t$ & Nodal & Galerkin\\
\midrule
$0.002$ & --- & $0.0024$\\
$0.004$ & $0.0009$ & $0.0033$\\
$0.008$ & $0.0015$ & $0.0026$\\
$0.020$ & $0.0232$ & $0.0198$\\
$0.040$ & $0.0131$ & $0.0131$\\
\bottomrule
\end{tabular}
\end{center}
Para $\Delta t\ge0.02$ ambos métodos pierden precisión de forma similar; para $\Delta t\le0.008$ el método nodal está más cerca de su referencia.

\paragraph{Experimento B1 ($t=0.04$, $L=6$).} Dependencia de la malla (nodal, $\Delta t=0.002$): la diferencia máxima entre $n=60$ y $n=150$ es $0.20$ (entre $n=100$ y $n=150$, $0.096$); con Galerkin, $0.20$ y $0.114$. El error espacial domina sobre las diferencias entre esquemas temporales ($\approx0.01$); las comparaciones de las secciones anteriores aíslan únicamente el error temporal a malla fija.

La diferencia máxima entre ambos métodos a igual $\Delta t$ es $0.012$ ($n=60$), $0.019$ ($n=100$) y $0.028$ ($n=150$) para $\Delta t=0.002$. En $n=150$ se varió $\Delta t$:
\begin{center}
\begin{tabular}{ccc}
\toprule
$\Delta t$ & $|{\rm nodal}-{\rm Galerkin}|$ & \\
\midrule
$0.004$ & $0.0239$ & \\
$0.002$ & $0.0279$ & \\
$0.001$ & $0.0230$ & \\
\bottomrule
\end{tabular}
\end{center}
y los cambios al reducir $\Delta t$ a la mitad fueron $0.0025$ y $0.0022$ (nodal) frente a $0.0148$ y $0.0148$ (Galerkin). Es decir, con esta implementación el método de Galerkin \emph{no converge} al reducir $\Delta t$ a malla fija (el resultado se desplaza $\approx0.015$ en cada reducción), mientras que el nodal sí se estabiliza. No se pudo determinar cuál de los dos se acerca más a la solución exacta del problema continuo, porque no se dispone de solución de referencia para B1. Posibles causas de la deriva del método de Galerkin, sin verificar: orden de cuadratura insuficiente para el integrando $\phi^n\circ X$ (no polinomial), y oscilaciones de la proyección $L^2$ cerca de la singularidad cónica de $\phi_0=r-r(\theta)$ en el origen, acumuladas en cada paso.

\subsection{Pérdida de la simetría radial con $r_0=1$}

Experimento: $A=0.5$, $G=20$, frontera inicial $r_0=1$ (solución exacta radial: $\dot R=G\,I_1(R)/I_0(R)-AGR/2$), dominio $[-4,4]^2$, método de características nodal con $\Delta t=0.002$, $t_{\max}=0.3$. Se mide la asimetría $(r_{\max}-r_{\min})/r_{\rm media}$ sobre 48 ángulos.

\begin{center}
\begin{tabular}{c|cc|cc|c}
\toprule
 & \multicolumn{2}{c|}{$n=60$} & \multicolumn{2}{c|}{$n=100$} & exacto\\
$t$ & $r_{\rm media}$ & asimetría & $r_{\rm media}$ & asimetría & $R(t)$\\
\midrule
$0.05$ & $1.2060$ & $0.0040$ & $1.2060$ & $0.0021$ & $1.2054$\\
$0.10$ & $1.4240$ & $0.0063$ & $1.4243$ & $0.0040$ & $1.4232$\\
$0.20$ & $1.8527$ & $0.0192$ & $1.8583$ & $0.0101$ & $1.8580$\\
$0.30$ & $2.2260$ & $0.0372$ & $2.2403$ & $0.0160$ & $2.2419$\\
\bottomrule
\end{tabular}
\end{center}

El radio medio es correcto (error $0.7\%$ con $n=60$, $0.07\%$ con $n=100$), pero la frontera se deforma de manera creciente con el tiempo: los radios en las direcciones de los ejes de la malla ($0^\circ,90^\circ,180^\circ,270^\circ$) crecen más que en las diagonales y entre ejes opuestos (p.\,ej.\ $270^\circ$ frente a $90^\circ$). La asimetría disminuye al refinar la malla ($0.037\to0.016$, razón $2.3$ para $h$ reducido en $1.67$). La malla con diagonales alternadas (\texttt{flags=1}) la reduce a $0.019$ con $n=60$, sin eliminarla.

\paragraph{Aislamiento del origen} (todas con $n=60$, $t=0.3$):
\begin{itemize}
	\item Con la curvatura exacta $\kappa=1/r$ en lugar de $\kappa_h$: asimetría $0.034$ (frente a $0.037$). \emph{La aproximación de la curvatura no es la causa.}
	\item Con la velocidad radial exacta $V(R(t))\,x/|x|$ en lugar de la calculada: asimetría $0.001$ y $r_{\rm media}=2.2412$. \emph{El transporte de $\phi$ conserva la simetría.}
	\item Diagnóstico de la velocidad calculada sobre el círculo $r=R_{\rm ex}$ ($t=0.12$, $R=1.51$): $\partial_nP$, que debería ser $0$ (en el caso radial $P$ es constante), varía entre $-0.12$ y $0.18$ con un patrón de periodo $90^\circ$ ligado a la malla ($-0.092$, $0.181$, $-0.093$, $0.071$ en $0^\circ,45^\circ,90^\circ,135^\circ$). $\partial_nU$ varía entre $0.588$ y $0.603$ (exacto $0.599$; media $0.596$), y al multiplicar por $G=20$ aporta una dispersión de $0.29$ en $V_n$. La velocidad calculada $v_D$ varía entre $4.17$ y $4.49$ (exacto $4.42$).
\end{itemize}
La causa es entonces el error, dependiente de la orientación de la malla, en el gradiente normal de las soluciones de $P$ y $U$ obtenidas con $\phi$-FEM, y no el transporte ni la curvatura. Ese error se integra en el tiempo y produce la deformación. Posibles remedios (no probados): calcular el flujo $\partial_n u$ de manera consistente (por ejemplo con la formulación dual con variable de flujo $y=-\nabla u$ que aparece comentada en el Capítulo 3, o recuperación de gradiente), usar grado mayor en $U$ y $P$, y emplear mallas simétricas.

\paragraph{Elementos de grado 3 (macro \texttt{gradPol3}).} Mismo experimento ($r_0=1$, $n=60$, $\Delta t=0.002$, $t_{\max}=0.3$), con $U$ y $P$ en $\mathbb{P}_3$, $\phi$ en $\mathbb{P}_4$, $\kappa$ en $\mathbb{P}_1$ y los mismos parámetros de estabilización ($\gamma=\sigma=20$). Asimetría y radio medio (exacto: $1.2054$, $1.4232$, $1.8580$, $2.2419$):

\begin{center}
\begin{tabular}{c|cc|cc|cc}
\toprule
 & \multicolumn{2}{c|}{$\mathbb{P}_2$ (original)} & \multicolumn{2}{c|}{$\mathbb{P}_3$, $S\in\mathbb{P}_1$} & \multicolumn{2}{c}{$\mathbb{P}_3$, $S\in\mathbb{P}_3$}\\
$t$ & $r_{\rm media}$ & asim. & $r_{\rm media}$ & asim. & $r_{\rm media}$ & asim.\\
\midrule
$0.05$ & $1.2060$ & $0.0040$ & $1.2057$ & $0.0028$ & $1.1326$ & $0.139$\\
$0.10$ & $1.4240$ & $0.0063$ & $1.4227$ & $0.0037$ & $1.1179$ & $0.385$\\
$0.20$ & $1.8527$ & $0.0192$ & $1.8529$ & $0.0100$ & $0.9682$ & $0.777$\\
$0.30$ & $2.2260$ & $0.0372$ & $2.2275$ & $0.0293$ & $0.7964$ & $0.912$\\
\bottomrule
\end{tabular}
\end{center}

\begin{itemize}
	\item En $t=0$ el $\mathbb{P}_3$ da $\partial_nU$ prácticamente exacto ($0.44638$ frente al exacto $0.44639$, dispersión $2\times10^{-5}$), pero la dispersión de $\partial_nP$ (que debe ser $0$) no mejora: $[-0.015,0.012]$ frente a $[-0.006,0.008]$ en $\mathbb{P}_2$.
	\item Con la extensión de la velocidad también en $\mathbb{P}_3$ ($S\in\mathbb{P}_3$) la simulación se desestabiliza (en $t=0.12$ las soluciones de $U$ y $P$ ya son incoherentes, p.\,ej.\ $\partial_nU\in[-0.07,0.72]$); en $t=0$ ya se observa $S\cdot n\in[3.56,3.90]$ frente a $v_D\approx3.93$.
	\item Con $S\in\mathbb{P}_1$ la asimetría mejora durante $t\le0.2$ (p.\,ej.\ $0.0100$ frente a $0.0192$ en $t=0.2$) pero crece rápido después ($0.0293$ en $t=0.3$, solo $21\%$ menos que en $\mathbb{P}_2$) y el radio medio no mejora. En $t=0.12$ aparecen valores atípicos locales ($\partial_nU=0.527$ en $45^\circ$ frente al exacto $0.599$; $v_D$ mínimo $2.08$).
\end{itemize}
Por tanto, con los parámetros actuales, pasar a $\mathbb{P}_3$ no resuelve la pérdida de simetría: el error pasa de una dispersión suave dependiente de la malla a valores atípicos puntuales. Es probable que $\gamma$ y $\sigma$ (elegidos para $\mathbb{P}_2$) deban aumentarse con el grado, pero eso no se ha probado.

\subsection{Curvatura por diferencias centrales en malla cuadrada}

Se implementó en \texttt{edp/crecimientotumoranocircularFD.edp} (variable \texttt{usekfd}) la curvatura
\[
\kappa=\frac{\phi_{xx}\phi_y^2-2\phi_x\phi_y\phi_{xy}+\phi_{yy}\phi_x^2}{|\nabla\phi|^3},
\]
con diferencias centrales de segundo orden de paso $d$ (lado de la malla cuadrada) evaluadas sobre $\phi_h$, y $\kappa_h$ interpolada en $\mathbb{P}_2$ en la banda. Se compara con la aproximación débil actual (Sección ``Aproximación de la curvatura'').

\paragraph{Exactitud estática} (error de $\kappa_h$ evaluada sobre la curva exacta, 360 puntos; sin errores relativos porque $\kappa$ se anula en puntos de inflexión):
\begin{center}
\begin{tabular}{lcccc}
\toprule
Caso & $n$ & Débil: error medio / máx. & Dif. centrales: error medio / máx.\\
\midrule
Círculo $R=2$ & $60$ & $1.5\times10^{-4}$ / $4.3\times10^{-4}$ & $3.1\times10^{-4}$ / $5.5\times10^{-4}$\\
Círculo $R=2$ & $100$ & $5.3\times10^{-5}$ / $1.6\times10^{-4}$ & $1.1\times10^{-4}$ / $2.0\times10^{-4}$\\
B1 & $60$ & $0.0146$ / $0.347$ & $0.0314$ / $0.344$\\
B1 & $100$ & $0.0045$ / $0.090$ & $0.0109$ / $0.100$\\
B1 & $150$ & $0.0021$ / $0.038$ & $0.0049$ / $0.044$\\
\bottomrule
\end{tabular}
\end{center}
Con paso $2d$ el error de B1 empeora ($0.109$, $0.042$, $0.019$ para $n=60,100,150$). En B1 la curvatura exacta varía entre $-2.9$ y $1.9$. Ambos métodos convergen aproximadamente con orden $h^2$, y las diferencias centrales tienen un error medio unas $2$ veces mayor que la aproximación débil con el mismo $h$. Con paso $d=h$ no son más exactas que el método actual.

\paragraph{Experimento $r_0=1$} ($n=60$, $\Delta t=0.002$, $t=0.3$): asimetría $0.030$ con diferencias centrales de paso $h$ ($0.034$ con paso $2h$), frente a $0.037$ con la aproximación débil y $0.034$ con la curvatura exacta $1/r$. La diferencia entre métodos de curvatura está dentro de lo que cambia al usar la curvatura exacta, por lo que la curvatura no explica la pérdida de simetría (ver la sección anterior).

\paragraph{Robustez en la banda.} Durante la evolución del círculo, $\max|\kappa_h-1/r|$ sobre todos los nodos de la banda fue $0.94$ con la aproximación débil en $t=0.3$ (con valores de $\kappa_h$ entre $-0.50$ y $1.03$, en nodos del borde de la banda) y $0.14$ con diferencias centrales (con valores entre $0.30$ y $0.60$). Las diferencias centrales no producen valores atípicos en el borde de la banda. En el experimento $r_0=1$ esto no cambió el resultado, pero no se probó en evoluciones largas de B1.

\subsection{Curvatura por la formulación variacional}

También se evaluó la formulación variacional del Capítulo 3 (bloque comentado): $N_h\in[\mathbb{P}_2]^2$ como proyección $L^2$ de $\nabla\phi_h/|\nabla\phi_h|$ sobre la banda y $K_h\in\mathbb{P}_2$ tal que
\[
\int_{\Omega_h^\Gamma}K_hv+\sigma_K h\sum_{E\in\mathcal{F}_h^\Gamma}\int_E[\partial_nK_h][\partial_nv]
=\int_{\partial\Omega_h^\Gamma}(N_h\cdot n)v-\int_{\Omega_h^\Gamma}N_h\cdot\nabla v ,
\]
con $\sigma_K=0.01$ (en B1 con $n=60$ el error medio fue $0.0198$, $0.0216$ y $0.0219$ para $\sigma_K=0.01,0.1,1$, es decir poco sensible a $\sigma_K$). Está disponible como \texttt{usekfd=2} en \texttt{edp/crecimientotumoranocircularFD.edp}.

Error medio / máximo de $\kappa_h$ sobre la curva exacta:
\begin{center}
\begin{tabular}{lccc}
\toprule
Caso & Débil (original) & Dif. centrales & Variacional\\
\midrule
Círculo $R=2$, $n=60$ & $1.5\times10^{-4}$ / $4.3\times10^{-4}$ & $3.1\times10^{-4}$ / $5.5\times10^{-4}$ & $1.4\times10^{-3}$ / $6.8\times10^{-3}$\\
Círculo $R=2$, $n=100$ & $5.3\times10^{-5}$ / $1.6\times10^{-4}$ & $1.1\times10^{-4}$ / $2.0\times10^{-4}$ & $7.8\times10^{-4}$ / $4.5\times10^{-3}$\\
B1, $n=100$ & $0.0045$ / $0.090$ & $0.0109$ / $0.100$ & $0.0073$ / $0.158$\\
B1, $n=150$ & $0.0021$ / $0.038$ & $0.0049$ / $0.044$ & $0.0040$ / $0.080$\\
\bottomrule
\end{tabular}
\end{center}
La formulación variacional es la menos exacta en el círculo (un orden de magnitud peor que la débil, con convergencia más lenta: $1.4\times10^{-3}\to7.8\times10^{-4}$ de $n=60$ a $100$) y queda entre las otras dos en el error medio de B1, con el mayor error máximo.

Experimento $r_0=1$ ($n=60$, $\Delta t=0.002$): asimetría $0.0276$ en $t=0.3$ (radio medio $2.2276$, exacto $2.2419$), dentro del rango de los demás métodos ($0.037$ débil, $0.030$ dif. centrales, $0.034$ curvatura exacta), de modo que tampoco explica la pérdida de simetría. En los nodos del borde de la banda $\max|\kappa_h-1/r|$ llegó a $1.10$ en $t=0.3$ (valores de $\kappa_h$ entre $0.28$ y $1.58$), peor que la aproximación débil ($0.94$) y que las diferencias centrales ($0.14$).
